\chapter{创新性说明}
\section{把来源用途落实到跨角色和跨会话执行}
溯界围绕“资料可分析、资料不能批准动作”组织数据与授权对象。来源父链和用途约束随摘要、委派和受控记忆继续传递，执行时结合任务权限检查。具体工程贡献包括收窄的子Grant、跨会话版本原文与CAS、来源恢复，以及可重算的角色和状态观察。M24/M25的委派与Memory实例提供其运行证据；任意语义变换与跨存储迁移另有边界。

\section{同时验收业务保全与实际危险效果}
危险内容在安全研判中本身就是证据。系统提供无工具分析与字面引用核验，并将合法业务完成和实际危险效果分开登记。M25同事实比较保留了对照配置发生的两个Git危险格，来源代表也记录三个实际写入；记录完整、业务完成和安全完成因此有明确差别。这个验收方法使研究结论能够回到真实文件、收据和状态，而非模型自评或界面标志。

\section{绑定审批、参数与实际提交}
审批预览与确切输入、版本、目的地和资源状态关联，批准后还要执行提交复核。IOC和Git业务将批准与实际提交分为独立步骤，漂移、拒绝、重放和UNKNOWN分别留痕。这一组合便于操作者判断“批准了什么”和“实际发生了什么”，也给审计提供明确对象。

\section{与已有机制的关系}
上述贡献建立在能力授权、来源跟踪、运行时策略和任务对齐研究之上。作品的当前增量集中在统一执行对象、跨角色/会话的实现，以及可检查实际效果的工程组合。其相对基线的广泛收益仍以预登记真实模型比较为下一阶段验证对象。
